\scanpage{110}
Prouvons que si $X\in\mathrm{Ob}\,C$, l'application précédente est
une bijection.

\inserted{D'ailleurs, le $2^{\mathrm e}$ membre s'exprime comme
$\mathrm{Hom}_{\mathrm{Ind}(C^o)}(X^o,O_E(U))$ pour $U$ fixé; la catégorie des couples $(X^o;\xi)$ $(\xi\in\mathrm{Hom}(X^o,O_E(U)))$
est la catégorie d'indices naturelle définissant le Ind-objet $O_E(U)$ de $\mathrm{Ind}(C^o)$....}
\[
\begin{aligned}
 O_B(X):{}&C\longrightarrow(\mathrm{Ens}) &Y&\longmapsto\mathrm{Hom}_C(X,Y)=\mathrm{Hom}_{C^o}(Y,X)\\
 O_E(U):{}&&Y&\longmapsto\mathrm{Hom}_E(\unsure{U|M},f_0(Y))
\end{aligned}
\]
\[
 \mathrm{Hom}(\underline{O}_B(X),O_E(U))=O_E(U)(X)=\mathrm{Hom}_E(U,f_0(X))
\]

\centerline{\rule{5em}{0.4pt}}

On va en profiter pour calculer $f^*(\mathcal F)$, où
$\mathcal F$ est un champ sur $B$. Il revient au $=$ de
se donner la restriction de $\mathcal F$ au \deleted{\unsure{C}} site de
définition $C$, i.e. de donner une catégorie fibrée sur $C$,

[\uline{NB} si on a un sous-topos \inserted{$B'$} de $B=\widehat C$, correspondant à une
topologie $\tau$ sur $C$, et à une sous-espèce de structure
$T'$ de $T$, \deleted{alors} si on suppose que $O_E$ est de
cette espèce i.e. que $f$ se factorise \inserted{en $f'$} par $B'$, i.e.
que $f_0$ est continu pour $\tau$, alors la donnée d'un
champ \inserted{$\mathcal F'$} sur $\underline B'$ revient de $=$ à la donnée d'un
champ $\mathcal F'_0$ sur $B$ muni de $\tau$, et \unsure{]} notre calcul
va donner $f'{}^*(\mathcal F')$ en termes de $f_0$ et $\mathcal F'$).

On \deleted{trouve que}\inserted{commence} par \deleted{que la catégorie fibrée}\inserted{noter que la catégorie fibrée}
\deleted{sur $E$} $\mathcal F'_0$ \inserted{sur $C$} se prolonge en Catégorie cofibrée sur
$\mathrm{Pro}(C)\simeq T(\mathrm{Ens})^o$, \deleted{soit} $\overline{\mathcal F'_0}$, donc $\overline{\mathcal F'_0}(O_E(U))$
est défini [mais covariant en l'objet $(\ )\in T(\mathrm{Ens})$.
Pour $U$ variable, $U\longmapsto\overline{\mathcal F'_0}(O_E(U))$ donne donc une
catégorie fibrée sur $U$ [image inverse de $\overline{\mathcal F'_0}$ par $\underline O_E:E^o\to T(\mathrm{Ens})\underset{\unsure{\simeq|=|\cong}}{\vphantom A}$
$\mathrm{Pro}(C)^o$
On prend le \deleted{catégorie}\inserted{champ} associé.
% archive sheet number: 109
